\section{Example 3: Loops with \acr{PHI} Nodes}

\textbf{TypeScript equivalent:}
\begin{lstlisting}[style=ts]
function sum(n: number): number {
  let total = 0;
  for (let i = 0; i < n; i++) {
    total += i;
  }
  return total;
}
\end{lstlisting}

\textbf{\acr{LLVM} \acr{IR} generation with loops:}
\begin{lstlisting}[style=ts]
import {
  LLVMCodegen,
  LLVMFunction,
  Parameter,
  BasicBlock,
  PhiInst,
  ICmpInst,
  BinaryInst,
  BrInst,
  RetInst
} from "@euriklis/llvm-ir";

const codegen = new LLVMCodegen("loop");

const sumFunc = new LLVMFunction({
  name: "sum",
  returnType: LLVMCodegen.types.i32
});

sumFunc.addParameter(new Parameter({
  type: LLVMCodegen.types.i32,
  name: "n"
}));

const entry = new BasicBlock("entry");
const loop = new BasicBlock("loop");
const exit = new BasicBlock("exit");

// entry: unconditional jump to loop
entry.setTerminator(new BrInst({ target: loop }));

// loop: PHI nodes for i and total
loop.addInstruction(new PhiInst({
  name: "i",
  type: LLVMCodegen.types.i32,
  incomingValues: [
    { value: 0, block: entry },
    { value: "i.next", block: loop }
  ]
}));

loop.addInstruction(new PhiInst({
  name: "total",
  type: LLVMCodegen.types.i32,
  incomingValues: [
    { value: 0, block: entry },
    { value: "total.next", block: loop }
  ]
}));

// total.next = total + i
loop.addInstruction(new BinaryInst({
  name: "total.next",
  opcode: LLVMCodegen.opcodes.add,
  type: LLVMCodegen.types.i32,
  lhs: "total",
  rhs: "i"
}));

// i.next = i + 1
loop.addInstruction(new BinaryInst({
  name: "i.next",
  opcode: LLVMCodegen.opcodes.add,
  type: LLVMCodegen.types.i32,
  lhs: "i",
  rhs: 1
}));

// cmp = (i.next < n)
loop.addInstruction(new ICmpInst({
  name: "cmp",
  predicate: LLVMCodegen.predicates.slt,
  type: LLVMCodegen.types.i32,
  lhs: "i.next",
  rhs: "n"
}));

// if cmp, continue loop; else exit
loop.setTerminator(new BrInst({
  condition: "cmp",
  conditionType: LLVMCodegen.types.i1,
  target: loop,
  falseTarget: exit
}));

// exit: return total
exit.setTerminator(new RetInst({
  type: LLVMCodegen.types.i32,
  value: "total"
}));

sumFunc.addBasicBlock(entry);
sumFunc.addBasicBlock(loop);
sumFunc.addBasicBlock(exit);
codegen.getModule().addFunction(sumFunc);

console.log(codegen.toString());
\end{lstlisting}

\textbf{Generated \acr{LLVM} \acr{IR}:}
\begin{lstlisting}[style=ir]
define i32 @sum(i32 %n) {
entry:
  br label %loop

loop:
  %i = phi i32 [ 0, %entry ], [ %i.next, %loop ]
  %total = phi i32 [ 0, %entry ], [ %total.next, %loop ]
  %total.next = add i32 %total, %i
  %i.next = add i32 %i, 1
  %cmp = icmp slt i32 %i.next, %n
  br i1 %cmp, label %loop, label %exit

exit:
  ret i32 %total
}
\end{lstlisting}

\textbf{Key concepts:}
\begin{itemize}[noitemsep]
  \item \texttt{PhiInst}: Merge values from different control flow paths
  \item Loop structure: \texttt{entry} → \texttt{loop} → (\texttt{loop} or \texttt{exit})
  \item \acr{SSA} form maintained: \texttt{i}, \texttt{i.next}, \texttt{total}, \texttt{total.next} are all unique
  \item Back-edge: \texttt{loop} can branch to itself
\end{itemize}

